\documentclass[12pt]{ltjsarticle}
\usepackage{geometry}
\geometry{margin=25mm}
\usepackage{amsmath,amssymb,amsthm,mathtools}
\usepackage{tikz-cd}
\usepackage{luatexja-fontspec}
\setmainfont{TeX Gyre Pagella}
\setsansfont{TeX Gyre Heros}
\setmonofont{TeX Gyre Cursor}

\title{逆行列の構造理解と形式検証メモ}
\author{TeX生成: CODEX (OpenAI) \\ Leanファイル生成: CODEX (OpenAI)}
\date{2025年9月19日}

\begin{document}

\maketitle

\section{概観}
Lean で作成した `zenla1\_session07.lean` の証明群をもとに、逆行列の基本的な性質を数式と図で整理する。特に $2\times2$、$3\times3$ 行列、そしてブロック行列表現に焦点を当て、Schur 補行列や転置との関係にも触れる。

\section{基本データ}
以下の行列を扱う。
\begin{align*}
A &= \begin{pmatrix} 3 & 2 \\ 1 & 4 \end{pmatrix}, &
B &= \begin{pmatrix} 2 & -1 \\ 3 & 2 \end{pmatrix}, &
C &= \begin{pmatrix} 1 & 2 & 3 \\ 0 & 2 & 1 \\ 0 & 0 & 3 \end{pmatrix}.
\end{align*}
$A$ の行列式は $\det(A)=3\cdot4-2\cdot1=10$ で、Lean 上では `norm_num` により即座に評価できる。

\section{$2\times2$ 行列の逆行列}
行列式が非零である $2\times2$ 行列
\[
M = \begin{pmatrix} a & b \\ c & d \end{pmatrix}
\]
に対し、逆行列は
\[
M^{-1} = \frac{1}{ad-bc}\begin{pmatrix} d & -b \\ -c & a \end{pmatrix}
\]
で与えられる。具体的に $A$ を代入すると
\[
A^{-1} = \frac{1}{10}\begin{pmatrix}4 & -2\\-1 & 3\end{pmatrix}
\]
となる。Lean の証明では成分ごとに `ext` で分解し、`simp` と `norm_num` で $AA^{-1}=I$ を示した。

\section{積・転置に関する性質}
$A,B$ の積 $AB$ は行列式 $70$ を持ち、可逆である。Lean では
\[
(AB)^{-1} = B^{-1}A^{-1}
\]
を各成分計算で検証した。さらに、
\[
(A^{\mathsf{T}})^{-1} = (A^{-1})^{\mathsf{T}}
\]
も成立し、転置と逆行列が可換であることが確認される。

\section{上三角行列の逆行列}
上三角行列 $C$ に対して、右逆行列
\[
C^{-1}_{\text{right}} = \begin{pmatrix}
1 & -1 & -\tfrac{2}{3}\\
0 & \tfrac{1}{2} & -\tfrac{1}{6}\\
0 & 0 & \tfrac{1}{3}
\end{pmatrix}
\]
を構成すると $CC^{-1}_{\text{right}} = I$ が成り立ち、逆行列も上三角になる。Lean では `simp` と `norm_num` を組み合わせて自動的に検証した。

\section{Schur 補行列とブロック行列}
$M$ を $1\times1$ ブロック 4 つからなる行列
\[
M = \begin{pmatrix} A_{11} & A_{12} \\ A_{21} & A_{22} \end{pmatrix}
\]
とする。$A_{11}$ が可逆で Schur 補行列 $S = A_{22} - A_{21}A_{11}^{-1}A_{12}$ が可逆ならば、
\[
M^{-1} = \begin{pmatrix}
A_{11}^{-1} + A_{11}^{-1} A_{12} S^{-1} A_{21} A_{11}^{-1} & -A_{11}^{-1} A_{12} S^{-1} \\
- S^{-1} A_{21} A_{11}^{-1} & S^{-1}
\end{pmatrix}
\]
が成立する。$A_{11}=3$, $A_{12}=2$, $A_{21}=1$, $A_{22}=4$ として評価すると $S = 10/3$、$(0,0)$ 成分は $2/5$ となり、Lean で求めた逆行列と一致する。

\begin{center}
\begin{tikzcd}
\mathbb{R}^2 \arrow[r, "A_{11}"] \arrow[d, "M"] & \mathbb{R}^2 \arrow[d, "A_{12}"] \\
\mathbb{R}^2 \arrow[r, "A_{21}"] & \mathbb{R}^2
\end{tikzcd}
\end{center}
この図式はブロック作用素の相互作用を表し、$A_{11}$ と $S$ の可逆性が全体の可逆性を保証することを視覚的に示している。

\section{Lean 証明の要点}
\begin{enumerate}
  \item 実数の逆数を扱うために `noncomputable section` を導入し、`Real` 上の非計算的構成を許容した。
  \item `change` を挿入してから `norm_num` を呼ぶことで、求める不等式を直接数値評価できる形に整えた。
  \item ブロック行列の逆行列は `simp` と `norm_num` のみで算出でき、Schur 補行列の公式と整合することを確認した。
\end{enumerate}

\section{まとめ}
Lean による形式的証明と数式計算を組み合わせることで、初等的な逆行列の性質を確実に裏付け、ブロック構造や Schur 補行列といった応用的テーマにも滑らかに拡張できることがわかった。今後は他のブロック分解（LU 分解や QR 分解）にも同様の手法を適用することで形式化の幅を広げられる。

\end{document}
